\begin{afterword}
\summaryhead

The post tells a parable. Suppose people had evolved the ability to add, and did not know how it worked. Calculators would be lookup tables typed in by experts, philosophers would say machines only simulate addition, and experts would forecast a century before a machine adds as well as a child. Researchers and bystanders would propose remedies: store more facts, learn like a child, read the Web, use evolution, let arithmetic emerge, wait for more computing power, scan a brain, train neural networks, invoke G\"odel's theorem or Searle.

The author draws two morals. First, beware of assertions you cannot regenerate from your own knowledge, as the calculators only play back sums they were given. Second, beware of dancing around a confusing gap instead of filling it; ``emergent'' and ``unknowable'' both avoid admitting that an insight is missing. The real history of AI, the author says, teaches the same. Pearl's graphical models explained what dozens of non-monotonic logics had patched case by case. The lesson: when the problem is your ignorance, clever ways around it fail, and ``Until you know your idea will work, it won't.''

\responsehead

The two morals are reasonable advice. A belief you cannot rebuild from your own understanding is fragile, and naming a gap ``emergent'' or ``unknowable'' does not fill it. The Pearl example is a real case of a problem that became easier once someone understood it.

The parable, though, contains its own answer. The reader knows that arithmetic is a simple formal system, so every voice in the list looks foolish by construction. Whether intelligence has a key of that kind is the question at issue, and a story in which the key exists cannot settle it. The post anticipates the charge of generalizing from fiction and says the lessons can be drawn from the real history of AI. It then describes one real insight, while its closing lesson is said to come from ``the history of previous key insights'' in AI.

The claim about academia also conflicts with the post's own example. The post says years of work on one problem is not a pursuit ``that academia is set up to permit.'' Pearl was at the University of California, Los Angeles, and Pearl's 1982 paper on the method acknowledges support from the National Science Foundation.

Later history bears on the rule, though the post could not have known it. One of the mocked voices proposed training neural networks to do arithmetic without knowing how they do it. Language models, neural networks trained on large amounts of text, did learn to add, and how they do it was worked out only afterwards, in 2025.

In short: a parable whose answer the reader already knows, two reasonable morals, and a lesson from AI's history that is based on one case.
\end{afterword}
